\section{GitHub Actions Workflow}

With the containers working locally, the next step was to automate the whole build and
push process using GitHub Actions. The idea is that nobody should have to run
\texttt{podman build} manually. Every time code is pushed to the main branch, the
pipeline should handle testing, building, scanning, and publishing the images on its own.

\subsection{Pipeline Structure}

The workflow is triggered on every push or pull request to main. I enabled concurrency
so that if a new push comes in while a run is still going, the old run gets cancelled.
Without this, two deployments can race against each other and produce inconsistent state.

The first job runs the backend tests. It checks out the repository, installs the right
version of Go with caching enabled, and runs \texttt{go test}. If the tests fail, the
rest of the pipeline does not run. This was an important thing to get right early,
because it is easy to build and push a broken image if you do not gate on tests first.

The second job builds the container images and pushes them to GHCR. I set up QEMU and
Docker Buildx so the images come out as multi-architecture, supporting both
\texttt{linux/amd64} and \texttt{linux/arm64}. Each image is tagged with the Git commit
SHA and the run number rather than \texttt{latest}, so every build is traceable.
Using \texttt{latest} for production images is a bad habit because you lose the ability
to know exactly what code is running. I configured the job to skip pushing on pull
requests, so images are only published when code is actually merged. I also added Trivy
scanning to catch critical and high-severity vulnerabilities before anything gets
deployed.

The third job updates the infrastructure repository with the new image tag. It checks
out the \texttt{guestbook-config} repo, uses \texttt{yq} to patch the image tag in the
deployment manifest, and pushes the change. This is the GitOps handoff point: the
application repo knows how to build, and the infrastructure repo records what version
is running. A GitHub token with write access to the config repo is required for this step.

The last job sends a build notification to Discord with the result, the version that was
built, and a link to the workflow run. This is a small addition but it makes it much
easier to know when something has deployed or broken without having to open the GitHub
UI.

\subsection{One Problem with Actions}

One error I ran into early was a confusing one. The workflow failed with:

\begin{lstlisting}
Error: Unable to resolve action actions/setup-buildx-action,
repository not found.
\end{lstlisting}

I had assumed all the standard actions were under the \texttt{actions/} namespace,
the same as \texttt{actions/checkout} and \texttt{actions/setup-go}. It turned out that
the Buildx and QEMU actions are maintained by Docker and live under the \texttt{docker/}
namespace. The fix was simply to change the paths to
\texttt{docker/setup-qemu-action} and \texttt{docker/setup-buildx-action}. It is the
kind of mistake you only make once.

\subsection{Two Repositories and Why}

At this point the project was split into two separate Git repositories. One holds the
application source code, Containerfiles, and the CI workflow. The other holds only the
Kubernetes and OpenShift manifest files. This separation follows a principle called
GitOps, where the desired state of the infrastructure is stored in Git and a tool like
ArgoCD keeps the cluster in sync with it.

The practical reason for having two repos is access control. A developer working on the
application code does not need cluster credentials. The CI workflow updates the config
repo with a new image tag, and ArgoCD picks that up and handles the deployment. There
is also a clean audit trail: if something breaks in production, you can look at the
commit history of the config repo and see exactly what changed and when.
